% 原书分章 PDF 第 13—17 页：第 16.4.2—16.4.3 节。
\subsection{行动机制 \texorpdfstring{$g_0^\theta$}{g0θ}}

在观察性研究中，实施上述 TMLE 和 IPW 估计量所需的条件概率 $g_0^\theta$ 未知，因而首先需要估计。下面先介绍如何基于五类倾向评分（propensity score，PS）分解这些干扰参数，再介绍本章用于估计这些 PS 的四种方法，包括两种基于模型的方法和两种数据自适应方法。

\paragraph{行动机制的分解。}
对于 $t=0,\ldots,t_0$，条件概率
\[
 g_0\bigl(A(t)\mid\bar L(t),\bar Y(t)=0,\bar A_1(t-1),\bar A_2(t-1)=0\bigr)
\]
称为时点 $t$ 的行动机制，它可基于以下五种 PS 作因子分解。
\begin{itemize}
\item TI 启动的 PS，记为 $\mu_1(t)$：
\[
 g_0\bigl(A_1(t)=1\mid\bar L(t),\bar Y(t)=0,
 \bar A_1(t-1)=0,\bar A_2(t)=0\bigr).
\]
\item TI 维持的 PS，记为 $\mu_2(t)$：
\[
 g_0\bigl(A_1(t)=1\mid\bar L(t),\bar Y(t)=0,
 \bar A_1(t-2),A_1(t-1)=1,\bar A_2(t)=0\bigr).
\]
\item 研究行政性结束所致右删失的 PS，记为 $\mu_3(t)$：
\[
 g_0\bigl(I\{A_2(t)=1,\Gamma=1\}=1\mid\bar L(t),\bar Y(t)=0,
 \bar A_1(t-1),\bar A_2(t-1)=0\bigr).
\]
\item 退出医疗保险计划所致右删失的 PS，记为 $\mu_4(t)$：
\[
\begin{aligned}
 g_0\bigl(I\{A_2(t)=1,\Gamma=2\}=1\mid&\bar L(t),\bar Y(t)=0,
 \bar A_1(t-1),\bar A_2(t-1)=0,\\
 &I\{A_2(t)=1,\Gamma=1\}=0\bigr).
\end{aligned}
\]
\item 死亡所致右删失的 PS，记为 $\mu_5(t)$：
\[
\begin{aligned}
 g_0\bigl(I\{A_2(t)=1,\Gamma=3\}=1\mid&\bar L(t),\bar Y(t)=0,
 \bar A_1(t-1),\bar A_2(t-1)=0,\\
 &I\{A_2(t)=1,\Gamma=1\}=0,
 I\{A_2(t)=1,\Gamma=2\}=0\bigr).
\end{aligned}
\]
\end{itemize}
因此，对于在时点 $t$ 之前未失效、且截至时点 $t$ 遵循决策规则 $d_\theta$ 的患者，干扰参数
\[
\begin{aligned}
 g_0\bigl(A(t)=d_\theta(t)(\bar L(t)\bar A(t-1))\mid
 \bar L(t),\bar Y(t)=0,\bar A(t-1)=d_\theta(\bar L(t-1))\bigr)
\end{aligned}
\]
的估计，可由上述五种 PS 的估计，按以下表达式导出。该表达式由时点 $t$ 行动机制的乘法法则分解得到：
\[
\begin{aligned}
 &\bigl(I\{\bar A_1(t-1)=0\}\mu_1(t)^{A_1(t)}
 \{1-\mu_1(t)\}^{1-A_1(t)}
 +I\{A_1(t-1)=1\}\mu_2(t)\bigr)\\
 &\hspace{5em}\times\{1-\mu_3(t)\}\{1-\mu_4(t)\}\{1-\mu_5(t)\}.
\end{aligned}
\]

\paragraph{按时间汇集数据的 logistic 模型。}
实践中，分析者常用的一种方法（Hernan 等，2000；Cole 等，2003），是使用跨时点 $t$ 汇集的数据拟合单一模型，称为“汇集模型”，从而同时估计所有时间区间中的上述每种 PS。具体而言，本研究汇集所有随访时点 $t$ 的数据，分别拟合仅含主效应项的 logistic 模型，估计三种右删失 PS，即 $\mu_3(t)$、$\mu_4(t)$、$\mu_5(t)$，以及 TI 维持 PS $\mu_2(t)$。我们还汇集所有 $t>0$ 时点的数据，拟合单一的仅主效应项 logistic 模型，估计 $t=0$ 之后的 TI 启动 PS，即 $t>0$ 时的 $\mu_1(t)$；另拟合一个单独的仅主效应项 logistic 模型，估计 $t=0$ 时的 TI 启动 PS $\mu_1(0)$。所谓“仅主效应项 logistic 模型”，是指对所考虑的每个解释变量仅包含其主效应项的 logistic 模型，即不含解释变量间的交互项。所考虑的解释变量包括全部不随时间变化的协变量，以及时变协变量的最近一次测量值（马尔可夫假设）。此外，估计三种右删失 PS 时，将上一期的 TI 暴露作为解释变量；估计所有 PS 时，均将 A1c 最近一次变化作为解释变量。所有汇集 logistic 模型还包括索引 90 天随访区间的变量，即 $t$。由此得到的干扰参数 $g_0^\theta$ 估计量记为 $g_n^\theta$。

\paragraph{按时间分层数据的 logistic 模型。}
在上一种方法中，每个汇集 logistic 模型都隐含假设：解释变量与 PS 结局变量的关联不随时间改变，例如 $\mu_5(t)$ 的结局变量是是否死亡。对这一假设的担忧（Platt 等，2009），促使我们改用不同的 logistic 模型，即“分层模型”，在每个时点 $t$ 分别估计每种 PS；或者至少在汇集模型中纳入解释变量与 $t$ 的函数之间的交互项。实际上，正是对混杂作用随时间改变的担忧，促使上一种方法指定两个不同模型来估计 TI 启动 PS：一个分层 logistic 模型估计 $\mu_1(0)$，另一个汇集 logistic 模型同时估计全部 $t>0$ 时的 $\mu_1(t)$。

为充分处理对混杂作用随时间改变的担忧，我们还考虑第二种 PS 估计方法：对每个时点 $t$，分别拟合五个仅主效应项 logistic 模型，估计五种 PS。这些分层模型的参数化形式与前一种方法相同，区别仅在于全部 logistic 模型均去掉时间变量 $t$。由此得到的干扰参数 $g_0^\theta$ 估计量记为 $g^\theta_{n,t}$。

\paragraph{数据自适应选择交互项的 logistic 模型。}
为减弱仅主效应项 logistic 模型施加的约束，我们考虑第三种 PS 估计方法：在上述分层 logistic 模型的基础上，以数据自适应方式纳入解释变量间的二阶交互项。本研究解释变量很多，又存在过拟合的顾虑，因此我们分别实施以下专门构造的数据自适应算法，从所有可能交互项中选择纳入每个分层模型的子集。对五种 PS 中的每一种以及每个时点 $t$，首先找出在单变量 logistic 回归中与 PS 结局变量关联最显著，即 $p$ 值最小的 15 个解释变量，据此计算 105 个二阶交互项。其次，对这 105 个交互项中的每一项，实施 PS 结局变量对该交互项及定义它的两个主效应项的 logistic 回归。第三，找出 $p$ 值低于 0.05 的交互项。最后，若满足这一标准的交互项超过 50 个，则只选取 $p$ 值最小的 50 项，将其加入该 PS 的分层 logistic 模型。所得干扰参数 $g_0^\theta$ 的估计量记为 $g^\theta_{n,t,\times}$。

\paragraph{超级学习。}
前述两种基于模型的估计量 $g_n^\theta$ 和 $g^\theta_{n,t}$，并不体现真实的领域知识，即所选 expit 函数是否适合正确表示这五种 PS 随时间变化的真实值。实际上，指定 PS 模型，例如选择 logistic 模型，通常源于传统、偏好或方便。为避免这种任意模型指定导致错误推断，实践中有人提出对干扰参数 $g_0^\theta$ 采用数据自适应估计；但此时一致估计又依赖于恰当地选择机器学习算法，亦称“学习器”。上一段考虑了这样一种学习器 $g^\theta_{n,t,\times}$，此外还有许多已提出的学习器，也可作为估计五种 PS 的候选方法。

与选择参数模型类似，选择学习器通常也不体现对不同可用学习器相对适用性的真实领域知识，因为“在实践中，通常不可能事先知道，对于给定的预测问题和数据集，哪种学习器表现最好”（van der Laan 等，2007）。为解决这一顾虑，可采用集成学习方法进行数据自适应估计，防范任意选择学习器造成的错误推断。超级学习（super learning，SL）就是这样一种方法，第 3 章已作详细讨论；它已知的理论性质支持其实际应用：就平均误差而言，超级学习器在渐近意义下的表现至少与所考虑的任何候选学习器一样好。如果超级学习器所考虑的某个候选学习器由一个参数模型定义，而且该模型恰好指定正确，那么，相比在正确模型下使用极大似然估计，使用 SL 的代价仅是预测变异性有限地增加。也可用其他集成学习方法替代 SL；不过，据我们所知，尚不能像超级学习那样，通过正式的有限样本和渐近结果，为这些替代方法的应用提供理论验证。

在这些二次分析中，我们对每个时点 $t$ 分别实施五个超级学习器，以估计五种 PS。每个超级学习器包含以下十个候选学习器：（a）五个由仅含主效应项的 logistic 模型定义的学习器，\footnote{通过 \texttt{SuperLearner} R 包中的 \texttt{SL.glm} 例程实现（Polley 和 van der Laan，2013）。}分别采用五个显著性水平 $\alpha=1\mathrm e{-30},1\mathrm e{-10},1\mathrm e{-5},0.1,1$，根据单变量回归中的显著 $p$ 值筛出最有预测力的解释变量；（b）五个多分类回归学习器，也采用相同的五个显著性水平，根据单变量回归中的显著 $p$ 值筛出最有预测力的解释变量。所得干扰参数 $g_0^\theta$ 的估计量记为 $g^\theta_{n,t,\mathrm{SL}}$。

\subsection{结局回归 \texorpdfstring{$Q_0^\theta$}{Q0θ}}

实施 TMLE 所需的条件期望 $Q_0^\theta$ 未知，因而首先需要估计。下面介绍作为干扰参数
\[
 E_{P_0}\bigl[Q^\theta_{L(k+2)}(\bar L(k+1))\mid
 \bar A(k)=d_\theta(\bar L(k)),\bar L(k),\bar Y(k)=0\bigr]
\]
初始估计量的两种数据自适应估计量。对于给定的 $k=t_0,\ldots,0$，仅使用在时点 $k$ 之前未失效、且截至 $k$ 遵循规则 $d_\theta$ 的患者数据，即 $\bar Y(k)=0$ 且 $\bar A(k)=d_\theta(\bar L(k))$ 的数据，来得到这些估计量。所得 TMLE 称为分层 TMLE。与之对应的汇集 TMLE，在构造 $Q_0^\theta$ 的初始估计量时，先汇集所有在时点 $k$ 之前未失效的患者数据，不论其先前是否遵循该动态干预，再将这些初始估计量在 $\bar A(k)=d_\theta(\bar L(k))$ 处求值。注意，在样本量较小的研究中，分层方法往往无法切实地为干扰参数 $Q_0^\theta$ 提供适当的初始估计；此时利用未经历相关处理历史的患者数据进行外推，可以改善 TMLE 的表现。

\paragraph{DSA。}
删除／替换／添加算法（Deletion/Substitution/Addition，DSA；Sinisi 和 van der Laan，2004；Neugebauer 和 Bullard，2010）实施基于交叉验证的数据自适应估计量选择过程，可作为估计条件期望的一种机器学习方法。本研究使用 DSA 构造 $Q_0^\theta$ 的初始估计量，候选估计量限制为仅含主效应项的 logistic 模型；候选解释变量包括全部不随时间变化的协变量、时变协变量的最近一次测量值，以及 A1c 最近一次变化。为缩短计算时间，我们只采用一次 5 折交叉验证划分，不执行删除和替换操作，并将最大模型规模，即 logistic 模型中的主效应项数量，设为 10。所得 $Q_0^\theta$ 的估计量记为 $Q^\theta_{n,\mathrm{DSA}}$。

\paragraph{超级学习器。}
基于与采用 SL 估计行动机制相同的理由，本研究也考虑用 SL 定义干扰参数 $Q_0^\theta$ 的初始估计量。实施细节见 Neugebauer 等（2014a）。具体而言，对每个时点 $k$，分别构造一个超级学习器，包含以下八类候选学习器；每个学习器使用不同的解释变量子集：（a）7 个由仅主效应项 logistic 模型定义的学习器；（b）5 个使用 AIC 逐步选择模型的学习器；（c）5 个神经网络学习器；（d）5 个 Bayes 回归学习器；（e）5 个多分类回归学习器；（f）5 个随机森林学习器；（g）5 个对分类树使用装袋法的学习器；（h）20 个使用平滑样条的广义可加模型学习器。考虑的解释变量包括 DSA 方法使用的全部变量，并扩展纳入这些变量间的二阶交互项，以及既往协变量的其他汇总量，如所有既往 A1c 测量的平均值，或既往 A1c 测量超过 8\% 的次数。之所以增加既往协变量的汇总量，是因为干扰参数 $Q^\theta_{L(k+1)}(\bar L(k))$ 按定义就是协变量历史 $\bar L(k)$ 的函数，如同每种 PS 都可能是既往观测协变量的函数一样。不过，实践中常可用马尔可夫假设，来论证仅依据最近的处理和协变量测量估计 PS 的合理性；对于估计干扰参数 $Q^\theta_{L(k+1)}(\bar L(k))$，为何采用类似方法是合理的，却并不清楚。所得 $Q_0^\theta$ 的估计量记为 $Q^\theta_{n,\mathrm{SL}}$。
